\section{MineDojo: convertir Minecraft en una Agent Research Platform abierta}

Con la recipe definida, el curso usa primero MineDojo para llevar los requisitos abstractos a una infraestructura concreta. MineDojo ofrece a la vez un environment API, un ecosistema de tareas y una internet knowledge database. Así, quien investiga no solo puede entrenar una policy, sino también estudiar la generación de tareas, el reward learning, la recuperación de conocimiento y las habilidades de largo plazo. Para entender esta sección hay que distinguir tres capas: qué puede simular el entorno, cómo se definen las tareas y qué tipo de supervisión pueden aportar los materiales de internet.

\subsection{Entorno, agente y base de conocimiento: un conjunto de tres piezas}

El diagrama del sistema de MineDojo divide el mundo abierto en tres objetos que se restringen entre sí: el Minecraft environment produce la información visual y el estado, el generalist agent elige la action y la internet knowledge base aporta demostraciones humanas y conocimiento previo en texto. No basta con entregarle páginas web al agent: hay que construir interfaces entre goal, observation, knowledge y control.

\lecturefigure{slide-08-minedojo-three-components.jpg}{Entorno abierto, agente y base de conocimiento de internet en MineDojo}{Diapositiva V008 recuperada de la grabación oficial; 00:10:06}

\begin{knowledgebox}{Lectura de la figura: seguir el flujo de datos, no los íconos}
Desde el entorno, el agente recibe la imagen en primera persona, el inventory y el estado; desde la base de conocimiento, obtiene materiales humanos como videos, Wiki y Reddit; a su vez, la action del agent modifica el entorno. La pregunta central es cómo se alinean ambos tipos de entrada con un mismo objetivo, no si los tres componentes son «todos muy grandes».
\end{knowledgebox}

\lecturefigure{slide-09-three-thousand-tasks-title.jpg}{MineDojo organiza la investigación de mundo abierto con más de 3,000 tareas}{Diapositiva V009 recuperada de la grabación oficial; 00:10:24}

«3,000+ tasks» es el punto de entrada a las tareas; no debe leerse como si la cantidad ya demostrara la generality. El valor de la escala de tareas está en cubrir distintos recursos, lugares, objetos, horizontes temporales y condiciones de éxito, de modo que se puedan observar la transfer y la compositionality. Si todas las tareas siguieran compartiendo un mismo patrón superficial, la cantidad por sí sola no produciría capacidad de mundo abierto.

\subsection{Programmatic Tasks y Creative Tasks}

El espacio de tareas de MineDojo contiene dos tipos de objetivos complementarios. Una programmatic task puede evaluarse automáticamente a partir del simulator state, por ejemplo, si el inventory contiene cierto objeto o si se derrotó a cierta criatura. Una creative task, en cambio, puede pedir «construir una casa bonita» o «apagar un incendio de forma creativa»; su criterio de éxito tiene componentes semánticos y estéticos, y no puede escribirse simplemente como un predicado de estado.

\lecturefigure{slide-10-task-categories.jpg}{Simulator, tareas programáticas y tareas creativas en MineDojo}{Diapositiva V010 recuperada de la grabación oficial; 00:10:48}

\begin{knowledgebox}{Términos clave: responsabilidades de los tres tipos de infraestructura}
\begin{tabularx}{\textwidth}{@{}L{0.20\textwidth}X X@{}}
\toprule
Componente & Problema que resuelve & Relación con el hilo del curso \\
\midrule
Versatile simulator & Un mundo observable, controlable y paralelizable & Aporta el ciclo cerrado activo y experimentos reproducibles \\
Programmatic task & Objetivos explícitos cuyo estado puede verificarse & Adecuada para reward automática y entrenamiento a gran escala \\
Creative task & Objetivos abiertos para los que es difícil escribir a mano un success predicate & Obliga al sistema a aprender una reward semántica y las preferencias humanas \\
\bottomrule
\end{tabularx}
\end{knowledgebox}

\lecturefigure{slide-11-inventory-tech-tree-combat.jpg}{Inventory, technology tree y combat forman un espacio de tareas combinatorio}{Diapositiva V011 recuperada de la grabación oficial; 00:11:18}

Esta diapositiva debe leerse primero por sus dependencias: obtener objetos avanzados suele requerir recolectar materias primas, fabricar herramientas, desbloquear recetas y enfrentar peligros; por eso, las tareas no son labels independientes, sino que comparten un mismo árbol tecnológico. Los objetos, herramientas y criaturas de la figura son solo ejemplos; lo que de verdad sustenta la transferencia es que las skills previas puedan combinarse en un nuevo plan de horizon más largo.

\lecturefigure{slide-12-creative-task-examples.jpg}{Las tareas creativas carecen de un criterio automático de éxito sencillo}{Diapositiva V012 recuperada de la grabación oficial; 00:11:57}

\begin{warningbox}{Creative no significa «asignar la reward de cualquier manera»}
Las tareas creativas son más difíciles de evaluar automáticamente, pero siguen requiriendo evidence explícita: el texto del objetivo, la semántica del video, las preferencias humanas, las restricciones del entorno o un learned reward model. Si se usa solo una puntuación ambigua, la policy puede encontrar un shortcut ajeno a la intención humana. Una evaluación abierta debe informar al mismo tiempo el origen de la reward, los casos de fallo y la revisión humana.
\end{warningbox}

\lecturefigure{slide-13-firefight-open-ended-demo.jpg}{La tarea abierta de apagar un incendio muestra el objetivo semántico y la secuencia de acciones}{Diapositiva V013 recuperada de la grabación oficial; 00:12:57}

Al ver la demo, primero hay que distinguir entre «la imagen parece razonable» y «la tarea se completó de forma verificable». El video muestra al agent frente al fuego, la estructura del entorno y la elección de herramientas, pero un único éxito no ofrece tasa de éxito, alcance de generalización ni reward robustness. Sobre todo ilustra que un objetivo abierto requiere conectar la semántica visual, el inventory y las acciones de varios pasos, en lugar de solo predecir el siguiente cuadro.

\teachervoice{El ponente distingue expresamente entre programmatic y creative tasks: en las primeras es fácil escribir el success criterion; en las segundas, incluso decidir «qué cuenta como completado» puede requerir un juicio semántico. Esta distinción explica directamente por qué después se introduce MineCLIP: la learned reward no es un módulo adicional, sino la interfaz por la que los objetivos abiertos entran en el RL.}

\subsection{Base de datos de internet: de la tabla de objetos a las preguntas humanas}

El entorno aporta experiencia, pero aprender desde cero, por trial-and-error, todas las recetas y estrategias resulta extremadamente costoso. Por eso MineDojo trata YouTube, Wiki y Reddit como un acervo de experiencia humana: los videos aportan conducta e información visual, Wiki aporta hechos estructurados y los foros aportan fallas, estrategias y problemas de cola larga. Las tres fuentes difieren en ruido y en dificultad de grounding, por lo que no pueden concatenarse simplemente como tokens homogéneos.

\lecturefigure{slide-14-internet-knowledge-sources.jpg}{YouTube, Wiki y Reddit constituyen fuentes de conocimiento complementarias}{Diapositiva V014 recuperada de la grabación oficial; 00:13:06}

\begin{knowledgebox}{Lectura de la figura: qué aporta cada una de las tres fuentes}
YouTube destaca en demostraciones dinámicas, pero no tiene action labels confiables; Wiki destaca en objetos, recetas y reglas, pero rara vez muestra la recuperación ante fallos; Reddit destaca en la resolución de problemas y la experiencia de cola larga, pero contiene ruido y juicios subjetivos. Un generalist agent necesita source-aware retrieval, en lugar de tratar todo el contenido como un mismo tipo de supervisión.
\end{knowledgebox}

\lecturefigure{slide-15-internet-knowledge-scale.jpg}{Materiales de internet a gran escala cubren el espacio visual y semántico de Minecraft}{Diapositiva V015 recuperada de la grabación oficial; 00:13:18}

Las miniaturas densas de esta diapositiva ilustran ante todo la coverage, no labels limpios. Al leer la figura, compare la variedad de objetos, escenas y puntos de vista; esto respalda que «internet contiene abundante conocimiento previo», pero no demuestra que los materiales ya estén alineados con el agent state, y mucho menos que todas las conductas de los videos sean correctas o reproducibles. La escala debe auditarse junto con la eliminación de duplicados, la calidad, los derechos de autor y la pertinencia para la tarea.

\lecturefigure{slide-16-minecraft-item-catalog.jpg}{Los catálogos de objetos, criaturas y acciones aportan conocimiento a nivel de entidad}{Diapositiva V016 recuperada de la grabación oficial; 00:14:15}

El catálogo de objetos ayuda al modelo a construir un entity vocabulary: qué objetos pueden recolectarse, combinarse, equiparse o resultan peligrosos. Primero conviene ver los límites entre categorías y después las posibilidades de acción que ofrece cada objeto respecto a cada action. Conocer el nombre no equivale a saber qué píxeles del punto de vista actual corresponden al objeto, ni a conocer las precondiciones de la acción; por eso siguen siendo necesarios el perception grounding y el environment feedback.

\lecturefigure{slide-17-crafting-knowledge.jpg}{El crafting knowledge conecta el objeto objetivo con sus materias primas y pasos}{Diapositiva V017 recuperada de la grabación oficial; 00:14:24}

Las recetas de fabricación convierten un declarative fact en una procedural constraint. Al leer la figura, rastree hacia atrás, desde el objeto objetivo, las materias primas, las herramientas y las condiciones de la mesa de trabajo; esta dependencia puede usarse para la plan decomposition y también para el curriculum. El diagrama respalda que «internet puede aportar pasos estructurados», pero durante la ejecución todavía hay que verificar el inventory, la ubicación y la recuperación ante fallos.

\lecturefigure{slide-18-potion-recipe-graph.jpg}{El grafo de recetas de pociones muestra cadenas largas de dependencias y planes combinados}{Diapositiva V018 recuperada de la grabación oficial; 00:14:33}

La dificultad de este tipo de grafo de dependencias no está en memorizar una recipe aislada, sino en mantener el estado a lo largo de varios objetos intermedios. Primero se identifica el nodo objetivo y luego se recorren las aristas hacia atrás para obtener cada ingredient necesario; después, el orden lógico del grafo se traduce en actions de navigation, collection y crafting dentro del mundo. Cualquier recurso faltante puede obligar al agent a replanificar.

\lecturefigure{slide-19-wiki-reddit-database.jpg}{Las páginas de Wiki y las discusiones de Reddit aportan reglas y experiencia de cola larga}{Diapositiva V019 recuperada de la grabación oficial; 00:14:42}

Las páginas web y los hilos de discusión contienen explicaciones humanas, alternativas y experiencias de fallos; son una teacher signal importante además de los videos. Al leer la figura, primero separe las páginas de hechos de las publicaciones de opinión y después revise la fecha de publicación, la versión y el contexto. Las respuestas del foro pueden inspirar estrategias, pero no deben ejecutarse sin verificación; el environment result es la evidencia final en la tarea actual.

\lecturefigure{slide-20-human-open-ended-questions.jpg}{Las preguntas abiertas planteadas por personas amplían la distribución de tareas}{Diapositiva V020 recuperada de la grabación oficial; 00:15:09}

Estas preguntas muestran que la distribución de objetivos proviene de jugadores reales y no de un conjunto limitado de plantillas diseñadas de antemano por los investigadores. Primero vea si la tarea requiere sentido común, planificación o creatividad, y después si puede calificarse automáticamente. Las preguntas sirven para construir benchmarks de task proposal y de retrieval, pero también exponen el problema de la evaluación: los objetivos en lenguaje pueden ser ambiguos y el éxito suele admitir varias realizaciones equivalentes.

\teachervoice{El ponente llama a los materiales de internet un reference manual: cuando una persona no sabe cómo hacer algo en Minecraft, también busca en Wiki, en videos y en foros. Lo esencial para el agent no es «memorizar todo internet», sino encontrar el conocimiento pertinente para un goal y un state concretos, y juzgar con los resultados de la ejecución si esa sugerencia realmente aplica.}

\lecturefigure{slide-21-generalist-agent-transition.jpg}{Del entorno y la base de conocimiento a un generalist agent que puede aprender}{Diapositiva V021 recuperada de la grabación oficial; 00:15:18}

Esta diapositiva de transición reduce el problema a lo siguiente: si ya se tienen un entorno abierto y conocimiento humano, ¿cómo se entrena un agent capaz de entender objetivos en lenguaje natural? La respuesta, en primer lugar, no es predecir acciones directamente, sino aprender un reward model capaz de alinear videos con objetivos en texto; ese modelo convierte objetivos semánticos difíciles de escribir a mano en retroalimentación que el RL puede usar.

\subsection{Resumen de la sección}

MineDojo reúne en una misma plataforma el entorno abierto, el ecosistema de tareas y el conocimiento de internet. Las programmatic tasks permiten el entrenamiento automático, mientras que las creative tasks revelan la carencia de evaluación semántica; YouTube, Wiki y Reddit aportan conocimiento complementario, pero este debe alinearse con el estado actual y verificarse mediante la ejecución. Esto establece la base de datos e interfaces para la learned reward de MineCLIP y para la línea de code/memory de Voyager.
